\documentclass[a4paper]{article}
\usepackage{setspace}
\doublespacing

\begin{document}
\section*{Optimisation and Parallelisation Approach}
\subsection*{Analysis of the Serial Implementation}
\begin{itemize}
    \item The serial implementation is based on Griebel's Nast2D software,
        provided in C and FORTRAN. For a given NACA airfoil body in a fixed-size
        problem space, the software iteratively solves for the two-dimensional
        velocity vectors and pressure scalars under a continuous west-to-east
        inward flow of incompressible fluid.
\end{itemize}
\subsection*{OpenMP Implementation Strategy}
\subsection*{Validation of Correctness}
\begin{itemize}
    \item To validate correctness, we visually inspect the VTK results produced
        by the serialiser. Automatic validation, such as using \texttt{diff}, is
        not useful here due to minor deviations in floating-point accuracy,
        compounded by the \emph{fast math} LLVM optimisations.
    \item In all cases but MPI, a single CPU thread produces the file. MPI makes
        use of the VTK's parallel writer capability in which each rank can
        produce independent files, then stitched together by metadata provided
        to the visualiser.
    \item Up to floating-point accuracy, all ports produce the same output for
        the test cases described here.
\end{itemize}
\section*{Performance Analysis}
\subsection*{Execution and Test Environment}
\subsection*{Serial Performance}
\subsection*{OpenMP Performance}
\subsection*{MPI Performance}
\subsection*{CUDA Performance}

\section*{Performance Comparison and Discussion}
\end{document}

